% Telos shared preamble.
%
% Every Telos publication is designed to be printed on a home laser printer,
% one sheet at a time, and carried into a boat or a garage. Two constraints
% follow from that and govern everything below:
%
%   1. Pure monochrome. No value is ever a hue. Emphasis is carried by rule
%      weight, fill, and type, never by color, so a grayscale or black-only
%      printer loses nothing.
%   2. Card density. A "one-pager" is a hard contract: the leaf must fit on a
%      single side of US Letter. The layout is therefore tight by default and
%      the environments below are sized to leave room for a plate.

\documentclass[10pt,letterpaper]{article}

\usepackage[margin=0.55in,includefoot,footskip=0.24in]{geometry}
\usepackage[T1]{fontenc}
\usepackage[utf8]{inputenc}
\usepackage{lmodern}
\usepackage{microtype}
\usepackage{array}
\usepackage{booktabs}
\usepackage{longtable}
\usepackage{tabularx}
\usepackage{enumitem}
\usepackage{needspace}
\usepackage{multicol}
\usepackage{ragged2e}
\usepackage{graphicx}
\usepackage{wrapfig}
\usepackage{xcolor}
\usepackage{fancyhdr}
\usepackage{titlesec}
\usepackage{hyperref}
\usepackage[most]{tcolorbox}
\usepackage{tikz}
\usetikzlibrary{arrows.meta,calc,positioning,shapes.geometric,decorations.pathmorphing,patterns,fadings}

% Semantic names are kept so a document reads intelligibly, but every value is
% black, white, or a neutral gray that survives a black-only print driver.
\definecolor{telosink}{gray}{0}
\definecolor{telospaper}{gray}{1}
\definecolor{telosrule}{gray}{0.35}
\definecolor{telosfill}{gray}{0.90}
\definecolor{telosdeep}{gray}{0.72}
\definecolor{telosmute}{gray}{0.45}

\hypersetup{hidelinks,pdfborder={0 0 0}}

% Keep an unchanged source byte-reproducible across rebuilds: no build-clock
% creation date and no random trailer ID.
\pdfinfoomitdate=1
\pdftrailerid{}

\setlength{\parindent}{0pt}
\setlength{\parskip}{0.42em}
\setlength{\tabcolsep}{4pt}
\setlength{\columnsep}{1.1em}
\renewcommand{\arraystretch}{1.18}
\setlist[itemize]{leftmargin=1.15em,itemsep=0.12em,topsep=0.2em,parsep=0pt}
\setlist[enumerate]{leftmargin=1.5em,itemsep=0.18em,topsep=0.2em,parsep=0pt}
\setlist[description]{leftmargin=0pt,itemsep=0.16em,topsep=0.2em,parsep=0pt,
  font=\normalfont\bfseries}

% ---------------------------------------------------------------------------
% Document identity. Each leaf declares these before \begin{document} so the
% running head, the colophon, and the PDF metadata all agree.
% ---------------------------------------------------------------------------
\newcommand{\TelosProject}{Telos}
\newcommand{\TelosDocTitle}{}
\newcommand{\TelosDocSubtitle}{}
\newcommand{\TelosRevision}{}

\newcommand{\telosmeta}[4]{%
  \renewcommand{\TelosProject}{#1}%
  \renewcommand{\TelosDocTitle}{#2}%
  \renewcommand{\TelosDocSubtitle}{#3}%
  \renewcommand{\TelosRevision}{#4}%
  \hypersetup{pdftitle={#2},pdfsubject={#3; version #4},
    pdfkeywords={Telos, version #4},pdfauthor={Telos}}%
}

\pagestyle{fancy}
\fancyhf{}
\renewcommand{\headrulewidth}{0pt}
\renewcommand{\footrulewidth}{0.4pt}
\fancyfoot[L]{\footnotesize\textsc{\TelosProject}}
\fancyfoot[C]{\footnotesize\TelosDocTitle}
\fancyfoot[R]{\footnotesize\thepage}

% The first page of a one-pager carries the masthead instead of a running foot,
% so its footer is the compact provenance line.
\fancypagestyle{telosfirst}{%
  \fancyhf{}%
  \renewcommand{\headrulewidth}{0pt}%
  \renewcommand{\footrulewidth}{0.4pt}%
  \fancyfoot[L]{\footnotesize\textsc{\TelosProject}}%
  \fancyfoot[C]{\footnotesize\TelosRevision}%
  \fancyfoot[R]{\footnotesize\thepage}%
}

% ---------------------------------------------------------------------------
% Masthead. A heavy rule, the title, a light subtitle, a thin rule. The
% optional argument is a right-aligned tag (a species' scientific name, a
% lesson number) that sits on the title's baseline.
% ---------------------------------------------------------------------------
\newcommand{\telosmasthead}[1][]{%
  \thispagestyle{telosfirst}%
  \begingroup
  \setlength{\parskip}{0pt}%
  \noindent\rule{\linewidth}{1.6pt}\par
  \vspace{0.30em}
  \noindent
  \begin{minipage}[b]{0.66\linewidth}
    \raggedright{\LARGE\bfseries\TelosDocTitle}
  \end{minipage}%
  \hfill
  \begin{minipage}[b]{0.32\linewidth}
    \raggedleft{\small\itshape #1}
  \end{minipage}\par
  \vspace{0.18em}
  \noindent{\small\TelosDocSubtitle}\par
  \vspace{0.28em}
  \noindent\rule{\linewidth}{0.4pt}\par
  \endgroup
  \vspace{0.35em}
}

% ---------------------------------------------------------------------------
% Headings. Compact, small-caps, rule-anchored — a card has no room for the
% article class's default vertical air.
% ---------------------------------------------------------------------------
\titleformat{\section}{\normalfont\large\bfseries}{}{0pt}{}[\vspace{-0.55em}\rule{\linewidth}{0.4pt}]
\titlespacing*{\section}{0pt}{0.85em}{0.35em}
\titleformat{\subsection}{\normalfont\normalsize\bfseries}{}{0pt}{}
\titlespacing*{\subsection}{0pt}{0.6em}{0.2em}
\titleformat{\subsubsection}{\normalfont\small\bfseries\scshape}{}{0pt}{}
\titlespacing*{\subsubsection}{0pt}{0.45em}{0.12em}

% A short run-in heading for card interiors that must not consume a full line.
\newcommand{\telosrubric}[1]{\textbf{#1}\enspace}

% Cross-reference helper. Every Telos project is a set of loose printed sheets,
% so a reference names the sheet rather than a page number.
\newcommand{\sheetref}[1]{\textit{see the \textbf{#1} sheet}}

% Which decision, ADR or sheet governs the paragraph you are reading. A reader
% who disagrees with an instruction should be able to find the reasoning behind
% it rather than having to take it on trust.
\newcommand{\governedby}[1]{{\footnotesize\textsc{governed by #1}}}

% ---------------------------------------------------------------------------
% Quantities. siunitx is not in this TeX Live split, and the guides need only a
% fixed thin space and upright units, so define that directly rather than take
% a package dependency the documented install would not satisfy.
%   \qty{9}{kV}  ->  9 kV        \unitof{kV}  ->  kV
% ---------------------------------------------------------------------------
\newcommand{\unitof}[1]{\ensuremath{\mathrm{#1}}}
\newcommand{\qty}[2]{#1\,\unitof{#2}}
\newcommand{\numrange}[2]{#1\,--\,#2}
\newcommand{\qtyrange}[3]{#1\,--\,#2\,\unitof{#3}}
% Inches and feet appear constantly in the fishing guide; keep one spelling.
\newcommand{\inch}[1]{#1\,in}
\newcommand{\feet}[1]{#1\,ft}

% ---------------------------------------------------------------------------
% Cards and boxes.
% ---------------------------------------------------------------------------

% Titled card with a hairline frame. The default card for facts a reader looks
% up rather than reads through.
\newtcolorbox{teloscard}[1]{
  enhanced, breakable,
  colback=telospaper, colframe=telosink, colbacktitle=telospaper,
  coltitle=telosink,
  boxrule=0.5pt, sharp corners,
  left=6pt, right=6pt, top=4pt, bottom=4pt,
  toptitle=3pt, bottomtitle=3pt, titlerule=0.4pt,
  before skip=6pt, after skip=6pt,
  title={#1}, fonttitle=\small\bfseries
}

% Untitled left-ruled block: hierarchy without a redundant label.
\newtcolorbox{telosframe}{
  enhanced, breakable,
  colback=telospaper, frame hidden, boxrule=0pt, sharp corners,
  borderline west={1.3pt}{0pt}{telosink},
  left=8pt, right=2pt, top=3pt, bottom=3pt,
  before skip=6pt, after skip=6pt
}

% Filled emphasis panel — the "read this first" block. The 10% fill stays
% legible under a cheap laser toner curve and still reads as distinct.
\newtcolorbox{telospanel}[1]{
  enhanced, breakable,
  colback=telosfill, colframe=telosink, colbacktitle=telosfill,
  coltitle=telosink,
  boxrule=0.5pt, sharp corners,
  left=6pt, right=6pt, top=4pt, bottom=4pt,
  toptitle=3pt, bottomtitle=3pt, titlerule=0.4pt,
  before skip=6pt, after skip=6pt,
  title={#1}, fonttitle=\small\bfseries
}

% Hazard block. Deliberately the heaviest object on any page: a 1.4pt frame
% plus a solid title bar. Nothing else in Telos is allowed to look like this,
% so a reader skimming a lesson cannot mistake it for ordinary emphasis.
\newtcolorbox{teloshazard}[1]{
  enhanced, breakable,
  colback=telospaper, colframe=telosink,
  colbacktitle=telosink, coltitle=telospaper,
  boxrule=1.4pt, sharp corners,
  left=6pt, right=6pt, top=5pt, bottom=5pt,
  toptitle=3pt, bottomtitle=3pt,
  before skip=8pt, after skip=8pt,
  title={\MakeUppercase{#1}}, fonttitle=\small\bfseries
}

% A stop-and-check gate. Used before anything destructive, and before anything
% that would be expensive to discover later. Shared by every project that has
% a procedure someone could get hurt following carelessly.
\newtcolorbox{gate}[1]{
  enhanced, breakable,
  colback=telospaper, colframe=telosink,
  boxrule=1.0pt, sharp corners,
  left=6pt, right=6pt, top=4pt, bottom=4pt,
  toptitle=3pt, bottomtitle=3pt, titlerule=0.4pt,
  before skip=7pt, after skip=7pt,
  colbacktitle=telospaper, coltitle=telosink,
  title={\MakeUppercase{#1}}, fonttitle=\small\bfseries
}

% ---------------------------------------------------------------------------
% Rating dots. A five-step scale that prints identically in black-only: filled
% versus open circles, never a shaded bar.
% ---------------------------------------------------------------------------
\newcommand{\telosdot}{\tikz[baseline=-0.42ex]\fill[telosink] (0,0) circle (0.28ex);}
\newcommand{\telosopendot}{\tikz[baseline=-0.42ex]\draw[telosink,line width=0.28pt] (0,0) circle (0.28ex);}
\makeatletter
\newcounter{telos@rate}
\newcommand{\telosrating}[1]{%
  \begingroup
  \setcounter{telos@rate}{0}%
  \@whilenum\value{telos@rate}<#1\do{\telosdot\kern0.22em\stepcounter{telos@rate}}%
  \@whilenum\value{telos@rate}<5\do{\telosopendot\kern0.22em\stepcounter{telos@rate}}%
  \endgroup
}
\makeatother

% ---------------------------------------------------------------------------
% Tables.
%
% tabularx reads its body by scanning for a literal \end{tabularx}, so it can
% never be wrapped in a \newenvironment. Documents therefore open tabularx and
% longtable directly and take their house style from these column types and the
% booktabs rules. The one exception is the fact strip below, which is built on
% plain tabular precisely so it *can* be an environment.
% ---------------------------------------------------------------------------
\newcolumntype{L}{>{\raggedright\arraybackslash}X}
\newcolumntype{R}{>{\raggedleft\arraybackslash}X}
\newcolumntype{C}{>{\centering\arraybackslash}X}
\newcolumntype{B}{>{\bfseries\raggedright\arraybackslash}X}
\newcolumntype{N}{>{\raggedright\arraybackslash}p{0.16\linewidth}}

% Key/value fact strip. Two flush columns of short lookups; the label column is
% a fixed fraction so a stack of cards aligns across the whole guide.
\newenvironment{telosfacts}[1][0.24]
  {\begingroup\small
   \setlength{\parskip}{0pt}%
   \renewcommand{\arraystretch}{1.25}%
   \begin{tabular}{@{}>{\bfseries\raggedright\arraybackslash}p{\dimexpr#1\linewidth\relax}@{\hspace{0.7em}}>{\raggedright\arraybackslash}p{\dimexpr\linewidth-#1\linewidth-0.7em\relax}@{}}}
  {\end{tabular}\par\endgroup}
\newcommand{\telosfact}[2]{#1 & #2\\}

% ---------------------------------------------------------------------------
% Seasonal / diurnal activity bar.
%
% \telosseasonbar{4,3,2,1,1,2,3,3,4,5,4,3} draws a twelve-cell month strip
% where each value 0-5 is a fill height. Height, not gray value, carries the
% signal, so it is exact under any toner curve; the cell is outlined either way
% so an empty month is still visibly a month.
% ---------------------------------------------------------------------------
\newcommand{\telosseasonlabels}{J,F,M,A,M,J,J,A,S,O,N,D}
\newcommand{\telosseasonbar}[1]{%
  \begingroup
  \begin{tikzpicture}[x=1.25em,y=2.1ex,baseline=0pt]
    \foreach \v [count=\i from 0] in {#1} {
      \draw[telosrule,line width=0.3pt] (\i,0) rectangle (\i+0.86,5);
      \ifnum\v>0
        \fill[telosink] (\i,0) rectangle (\i+0.86,\v);
      \fi
    }
    \foreach \m [count=\i from 0] in \telosseasonlabels {
      \node[font=\tiny,anchor=north,inner sep=1.2pt] at (\i+0.43,-0.1) {\m};
    }
  \end{tikzpicture}%
  \endgroup
}

% Four-cell day-part strip: dawn, midday, dusk, night.
\newcommand{\telosdaylabels}{Dawn,Mid,Dusk,Night}
\newcommand{\telosdaybar}[1]{%
  \begingroup
  \begin{tikzpicture}[x=2.9em,y=2.1ex,baseline=0pt]
    \foreach \v [count=\i from 0] in {#1} {
      \draw[telosrule,line width=0.3pt] (\i,0) rectangle (\i+0.92,5);
      \ifnum\v>0
        \fill[telosink] (\i,0) rectangle (\i+0.92,\v);
      \fi
    }
    \foreach \m [count=\i from 0] in \telosdaylabels {
      \node[font=\tiny,anchor=north,inner sep=1.2pt] at (\i+0.46,-0.1) {\m};
    }
  \end{tikzpicture}%
  \endgroup
}

% A bar needs vertical room for its labels when it sits in a fact strip.
\newcommand{\telosbarrow}[1]{\rule{0pt}{4.2ex}#1\rule[-1.6ex]{0pt}{1.6ex}}

% ---------------------------------------------------------------------------
% Plates. A species or apparatus illustration with a caption rule beneath it.
% \telosplate{width fraction}{file}{caption}
% ---------------------------------------------------------------------------
\newcommand{\telosplate}[3]{%
  \begingroup
  \centering
  \includegraphics[width=#1\linewidth]{#2}\par
  \vspace{0.15em}
  \begin{minipage}{#1\linewidth}
    \rule{\linewidth}{0.4pt}\par
    \vspace{0.1em}
    {\footnotesize #3\par}
  \end{minipage}\par
  \endgroup
  \vspace{0.3em}
}

% TikZ drawing conventions shared by every rig, knot, and circuit diagram.
\tikzset{
  telos line/.style={draw=telosink,line width=0.7pt},
  telos hair/.style={draw=telosink,line width=0.35pt},
  telos heavy/.style={draw=telosink,line width=1.2pt},
  telos node/.style={draw=telosink,fill=telospaper,align=center,inner sep=4pt,font=\small},
  telos emphasis/.style={draw=telosink,line width=1.2pt,fill=telospaper,align=center,
    inner sep=4pt,font=\small\bfseries},
  telos label/.style={font=\scriptsize,inner sep=1.5pt},
  telos arrow/.style={-{Stealth[length=1.8mm]},draw=telosink,line width=0.6pt},
  telos dim/.style={|<->|,draw=telosink,line width=0.35pt},
  telos water/.style={draw=telosrule,line width=0.4pt,decorate,
    decoration={snake,amplitude=0.4mm,segment length=3mm}},
}

% ---------------------------------------------------------------------------
% Lesson furniture
% ---------------------------------------------------------------------------

% \lessonhead{number}{time}{who does what}
\newcommand{\lessonhead}[3]{%
  \par\noindent
  \begin{minipage}{\linewidth}
    \rule{\linewidth}{0.4pt}\par
    \vspace{0.3em}
    {\small\textbf{Lesson #1}\quad$\cdot$\quad #2\quad$\cdot$\quad #3\par}
    \vspace{0.25em}
    \rule{\linewidth}{0.4pt}
  \end{minipage}\par
  \vspace{0.5em}}

% The materials list. Two columns of short lines; a lesson that needs a long
% shopping list is a lesson that has not been thought through.
\newenvironment{materials}
  {\begin{telospanel}{What you need}\begin{multicols}{2}
   \begin{itemize}[leftmargin=1em,itemsep=0.05em,topsep=0pt]}
  {\end{itemize}\end{multicols}\end{telospanel}}

% The four rubrics every lesson carries, in the same order every time:
% what you are about to see, what to do, what it means, and what to ask.
\newcommand{\lessonsection}[1]{\section*{#1}\vspace{-0.35em}\rule{\linewidth}{0.4pt}\vspace{0.2em}}

% A prediction box. The point of the curriculum is that they commit to an
% answer before the demonstration, not after.
\newtcolorbox{predict}{
  enhanced, breakable,
  colback=telospaper, colframe=telosink,
  boxrule=0.5pt, sharp corners,
  left=6pt, right=6pt, top=4pt, bottom=4pt,
  before skip=6pt, after skip=6pt,
  borderline west={2.2pt}{0pt}{telosink},
  title={\small\bfseries Predict first}, fonttitle=\small\bfseries,
  colbacktitle=telospaper, coltitle=telosink,
  toptitle=3pt, bottomtitle=3pt, titlerule=0.4pt,
}

% ---------------------------------------------------------------------------
% Colophon. Rendered once, at the end of the last page, in the recovered footer
% area so it never creates a page of its own.
% ---------------------------------------------------------------------------
\newif\ifTelosColophonRendered
\newcommand{\TelosColophonExtra}{}
\newcommand{\TelosColophon}{%
  \ifTelosColophonRendered\else
  \global\TelosColophonRenderedtrue
  \par
  \enlargethispage{0.5in}%
  \vspace{0.3em}%
  \begingroup
  \setlength{\parskip}{0pt}%
  \begin{minipage}{\linewidth}
  \rule{\linewidth}{0.35pt}\par
  \vspace{0.3em}%
  \fontsize{7}{8.2}\selectfont
  \textbf{Telos.} \TelosProject{} — \TelosDocTitle. Revised \TelosRevision.
  \TelosColophonExtra{}
  Project-created text and diagrams are licensed \href{https://creativecommons.org/licenses/by/4.0/}{CC BY 4.0}.
  Incorporated public-domain artwork remains public domain and is credited on
  the plate it appears with. See \texttt{LICENSE} and \texttt{CREDITS.md} in the
  source. Nothing here is professional advice; verify current regulations and
  follow the stated safety procedure.
  \end{minipage}
  \endgroup
  \fi
}
\AtEndDocument{\TelosColophon}

% Shared macros for the homelab project.
%
% This project's documents are read in two very different situations: at a desk
% while designing, and standing in front of a broken machine at eleven at night
% with no network. The second case is the one that sets the style. Commands are
% set apart so they can be found by eye, every dangerous step has a visible
% verification after it, and nothing depends on being able to reach a web page.

% ---------------------------------------------------------------------------
% Decision-record entries, used by the generated digest.
% ---------------------------------------------------------------------------
\newcommand{\adrentry}[4]{%
  \par\vspace{0.7em}
  \needspace{4\baselineskip}%
  \noindent
  \begin{minipage}{\linewidth}
    \rule{\linewidth}{0.6pt}\par
    \vspace{0.25em}
    {\small\textbf{ADR #1}\quad #2\par}
    \vspace{0.1em}
    {\fontsize{7}{8.2}\selectfont\textsc{#3}\quad$\cdot$\quad #4\par}
    \vspace{0.2em}
    \rule{\linewidth}{0.3pt}
  \end{minipage}\par
  \vspace{0.3em}}

% ---------------------------------------------------------------------------
% Procedures.
%
% A rebuild manual is a sequence of commands and the observable evidence that
% each one worked. These environments keep that pairing structural rather than
% a matter of the author remembering.
% ---------------------------------------------------------------------------

% A block of shell to type. Deliberately not a listing package: plain verbatim
% in a ruled box prints identically everywhere and has no font dependency.
\newenvironment{shellblock}
  {\par\vspace{0.25em}\noindent
   \begin{minipage}{\linewidth}
   \rule{\linewidth}{0.3pt}\par\vspace{0.15em}
   \begingroup\footnotesize\ttfamily\obeylines\obeyspaces\parindent=0pt\parskip=0pt}
  {\endgroup\par\vspace{0.15em}\rule{\linewidth}{0.3pt}\end{minipage}\par\vspace{0.25em}}

% What you should see if the previous step worked. A step without one of these
% is a step you cannot verify, and in a rebuild manual that is a defect.
\newcommand{\expectthis}[1]{%
  \par\vspace{0.1em}
  \noindent\begin{minipage}{\linewidth}
  \hspace*{0.6em}\begin{minipage}{\dimexpr\linewidth-0.6em\relax}
  {\footnotesize\textbf{Expect:} #1\par}
  \end{minipage}\end{minipage}\par\vspace{0.3em}}

% \begin{gate} now lives in common/preamble.tex, shared with the other
% projects that document procedures.

% \governedby now lives in common/preamble.tex.

% An instance value that comes from the private overlay. In the published
% documents these render as visible placeholders, which is deliberate: a
% published manual must be complete and obviously generic, never look like it
% is describing somebody's actual network.
\newcommand{\instancevalue}[2]{\texttt{\textlangle #1\textrangle}%
  \footnote{Instance value. Supplied from \texttt{homelab/instance/} at build
  time; #2.}}
\newcommand{\placeholder}[1]{\texttt{\textlangle #1\textrangle}}

% ---------------------------------------------------------------------------
% Role and boundary diagrams.
% ---------------------------------------------------------------------------
\tikzset{
  role/.style={draw=telosink,line width=0.9pt,fill=telospaper,align=center,
    inner sep=5pt,font=\small,minimum width=2.6cm},
  roleoff/.style={draw=telosmute,line width=0.5pt,dash pattern=on 2.5pt off 1.8pt,
    fill=telospaper,align=center,inner sep=5pt,font=\small,minimum width=2.6cm},
  boundary/.style={draw=telosmute,line width=0.5pt,dash pattern=on 3pt off 2pt,
    rounded corners=3pt},
  flow/.style={-{Stealth[length=2mm]},draw=telosink,line width=0.7pt},
  hlabel/.style={font=\fontsize{6.4}{7.4}\selectfont},
}


\telosmeta{Homelab}{Workstation Factory}
  {From an isolated bootstrap VM to repeatable Windows 11 and Arch laptops}{20260727.001}

\tikzset{
  wf pencil/.style={draw=black!78,line width=.48pt,line cap=round,line join=round},
  wf faint/.style={draw=black!32,line width=.27pt,line cap=round,line join=round},
  wf flow/.style={-{Stealth[length=1.8mm]},draw=black!72,line width=.55pt},
  wf deny/.style={draw=black!62,line width=.45pt,dashed}
}
\newcommand{\wfcaption}[1]{\par\smallskip{\centering\footnotesize\itshape #1\par}}

\begin{document}
\telosmasthead[Operator manual]

This manual creates a dual-boot laptop that can authenticate the same person
under Windows 11 Pro or Arch Linux. It starts with an isolated Arch virtual
machine named \texttt{bootstrap-dc}; it does not require a permanent Controller.
The public procedure uses placeholders. Put household names, addresses,
machine inventory and encrypted secret material in the private instance
repository.

\begin{gate}{Phase-1 boundary}
UniFi remains the only DHCP authority. The Controller hosts Samba AD with
integrated DNS, and HTTP/iPXE services. Workstation disks are deliberately
\textbf{not encrypted}; Secure Boot is deliberately disabled. Cached domain
login may work indefinitely while a laptop is away, so a disconnected laptop
cannot be remotely revoked. Treat every phase-1 workstation as unsuitable for
unrecoverable or sensitive data.
\end{gate}

\section{The result}

\begin{center}
\begin{tikzpicture}[x=1cm,y=1cm,font=\scriptsize]
  % Gateway appliance.
  \draw[wf pencil,rounded corners=1pt] (-5.8,.05) rectangle (-3.65,1.55);
  \draw[wf faint] (-5.62,.23) rectangle (-3.83,1.37);
  \foreach \y in {.43,.69,.95}
    {\draw[wf pencil] (-5.42,\y) rectangle (-4.18,\y+.14);
     \fill[black!58] (-3.99,\y+.07) circle (.035);}
  \node[align=center] at (-4.72,1.10) {UniFi\\gateway};
  \node[align=center] at (-4.72,-.18) {addresses + policy};

  % Host and VM.
  \draw[wf pencil,rounded corners=1pt] (-1.9,-.05) rectangle (1.35,1.70);
  \draw[wf faint] (-1.72,.12) rectangle (1.17,1.53);
  \draw[wf pencil] (-1.42,.37) rectangle (.87,1.25);
  \draw[wf faint] (-1.27,.52) rectangle (.72,1.10);
  \node[align=center] at (-.28,.83) {\texttt{bootstrap-dc}\\Arch VM};
  \node[align=center] at (-.28,-.18) {AD/DNS + boot};

  % Laptop.
  \draw[wf pencil,rounded corners=1pt] (3.60,.34) rectangle (5.85,1.72);
  \draw[wf faint] (3.74,.48) rectangle (5.71,1.58);
  \draw[wf pencil] (3.25,.03)--(6.20,.03)--(5.85,.34)--(3.60,.34)--cycle;
  \draw[wf faint] (4.22,.20) rectangle (5.23,.29);
  \node[align=center] at (4.73,1.05) {Windows 11\\Arch Linux};
  \node[align=center] at (4.73,-.18) {one domain identity};

  \draw[wf flow] (-3.65,.87)--(-1.90,.87)
    node[midway,above,align=center]{DHCP\\DNS route};
  \draw[wf flow] (1.35,.87)--(3.60,.87)
    node[midway,above,align=center]{PXE + AD\\authentication};
\end{tikzpicture}
\wfcaption{Stable names and AD discovery decouple each workstation from the temporary VM.}
\end{center}

\begin{telosfacts}[0.22]
\telosfact{Identity}{The permanent DNS domain comes from \texttt{identity\_dns\_domain}. Clients discover domain controllers through DNS service records, never a baked-in machine address.}
\telosfact{Boot}{Windows is the default after a five-second UEFI boot menu; Arch remains independently bootable.}
\telosfact{Storage}{The default ratio is 75 percent Windows and 25 percent Arch after fixed partitions, subject to 160 GiB and 64 GiB minimums.}
\telosfact{Network home}{A local profile always opens. Optional SMB storage mounts afterward and fails soft.}
\end{telosfacts}

\section{Stage 0 --- Record inputs without secrets}

Create a private instance record before changing the network. It names the
site prefix, VLANs, subnets, gateway, Controller addresses, workstation
hostname, disk serial, desired layout and physical switch ports. It contains
no plaintext password, Wi-Fi key, recovery material or reusable join token.

\begin{tabularx}{\linewidth}{@{}p{.25\linewidth}p{.33\linewidth}L@{}}
\toprule
\textbf{Public constant} & \textbf{Private instance value} & \textbf{Secret channel}\\
\midrule
\texttt{identity\_dns\_domain}; Windows 11 Pro; UEFI & Addresses, hostnames, serials, MACs, port assignments & AD bootstrap password, join credential, Wi-Fi key\\
75/25 default; 160/64 GiB floors & Approved per-machine override and disk identity & Local rescue and user passwords\\
\texttt{boot\_service\_fqdn} alias & Current server behind the alias & NAS and deployment credentials\\
\bottomrule
\end{tabularx}

\askbefore{Does the machine name satisfy both DNS and Windows naming rules, and
is it absent from AD and DNS? Does the disk's physical label match the recorded
model, serial and capacity? Who owns the switch port and DHCP scope?}

\section{Stage 1 --- Build the isolated bootstrap VM}

Run the repository dependency target on an Arch build host, then create the
virtual machine. The VM has four vCPUs, 8 GiB RAM and an 80 GiB disk. During
this stage it attaches only to an isolated virtual network.

\begin{shellblock}
make homelab-bootstrap-deps
make homelab-bootstrap-vm-create APPLY=1
make homelab-bootstrap-vm-run APPLY=1
make homelab-bootstrap-controller INVENTORY=<private-inventory>
\end{shellblock}

\expectthis{The VM reports hostname
\texttt{bootstrap\_dc\_fqdn}; Samba, Kerberos and DNS tests pass; no
packet from the VM reaches the household LAN.}
\failurebranch{If a target is missing, reports a proposed command instead of
running it, or cannot prove isolation, stop. Do not improvise package commands
or bridge the VM manually; fix the reproducible target.}

\begin{center}
\begin{tikzpicture}[x=1cm,y=1cm,font=\scriptsize]
  \draw[wf pencil] (-5.55,-.25) rectangle (-2.55,1.38);
  \draw[wf faint] (-5.37,-.07) rectangle (-2.73,1.20);
  \node[align=center] at (-4.05,.71) {Arch build host\\Make + QEMU};
  \draw[wf pencil] (-1.00,-.03) rectangle (1.55,1.16);
  \draw[wf faint] (-.84,.13) rectangle (1.39,1.00);
  \node[align=center] at (.28,.58) {\texttt{bootstrap-dc}\\isolated NIC};
  \draw[wf flow] (-2.55,.57)--(-1.00,.57);
  \draw[wf deny] (2.45,-.15)--(2.45,1.25);
  \draw[wf pencil] (2.20,1.00)--(2.70,.10);
  \draw[wf pencil] (2.20,.10)--(2.70,1.00);
  \draw[wf faint] (3.35,-.25) rectangle (5.65,1.38);
  \node[align=center] at (4.50,.62) {household LAN\\unchanged};
\end{tikzpicture}
\wfcaption{The first proof is negative: the unfinished services cannot reach the live LAN.}
\end{center}

\proofpair{make homelab-test}{The repository's VM, identity, PXE and
workstation tests pass; save the test report and artifact digests.}
\proofpair{make homelab-bootstrap-vm-status}{Reports the expected isolated VM
state without changing the VM or UniFi.}

\section{Stage 2 --- Establish the permanent identity plane}

The first domain controller creates the permanent forest
\texttt{identity\_dns\_domain}, \texttt{kerberos\_realm}, and
\texttt{netbios\_name}. The Controller's implementation may later move to a VM or
physical machine; those names do not change.

\begin{gate}{DNS is part of Active Directory}
Windows and Arch locate LDAP and Kerberos through DNS SRV records. All
domain-member workstation interfaces must query AD DNS. The AD DNS service may
forward unrelated queries, but a public resolver cannot replace it.
\end{gate}

\proofpair{samba-tool domain level show}{The forest and domain levels are
reported without error.}
\proofpair{host -t SRV \_ldap.\_tcp.\placeholder{identity-domain}}{At least one domain
controller target and port 389 are returned.}
\proofpair{host -t SRV \_kerberos.\_udp.\placeholder{domain}}{At least one domain
controller target and port 88 are returned.}
\proofpair{kinit <domain-admin> \&\& klist}{A ticket for
\placeholder{kerberos-realm} exists; destroy it after the test with \texttt{kdestroy}.}

\section{Stage 3 --- Build versioned boot targets}

Build each payload independently, then build the signed-off set. The Windows
target consumes a locally supplied Microsoft ISO and never redistributes it.
Every release directory is immutable and contains a SHA-256 manifest.

\begin{shellblock}
make homelab-pxe-controller
make homelab-pxe-arch
make homelab-pxe-windows WINDOWS\_ISO=/private/path/windows-11.iso
make homelab-pxe-all
make homelab-pxe-test
\end{shellblock}

\expectthis{The test boots each menu entry in an isolated UEFI virtual machine,
fetches every payload over HTTP, verifies each digest, and stops before disk
writes unless a disposable test disk and full serial authorization are
present.}
\failurebranch{A missing Microsoft driver, changed ISO digest or failed fetch
invalidates the release. Never edit an artifact in place; repair the source and
build a new version.}

\begin{center}
\begin{tikzpicture}[node distance=8mm and 8mm,font=\scriptsize]
  \node[role] (source) {source\\+ local ISO};
  \node[role,right=of source] (build) {build\\isolated};
  \node[role,right=of build] (test) {UEFI test\\all targets};
  \node[role,right=of test] (release) {release\\manifest};
  \node[role,right=of release] (current) {\texttt{current}\\atomic link};
  \draw[flow] (source)--(build);
  \draw[flow] (build)--(test);
  \draw[flow] (test)--(release);
  \draw[flow] (release)--(current);
  \draw[flow,dashed] (test.south) to[bend right=30]
    node[hlabel,below]{failure: new build}(source.south);
\end{tikzpicture}
\wfcaption{Promotion changes one pointer; rollback selects the previous verified release.}
\end{center}

\section{Stage 4 --- Choose the real network attachment}

This is a deployment gate, not a build prerequisite. Choose exactly one:
existing LAN, a dedicated second interface, or a VLAN trunk. Record the reason,
the recovery path and the observed switch-port configuration. UniFi remains
the sole address authority in every choice.

\begin{gate}{Recommended production boundary}
Use a dedicated provisioning VLAN with UniFi DHCP options pointing UEFI
clients at the stable boot endpoint. Permit only AD/DNS, boot HTTP, time,
approved package sources and explicitly tested storage. A final deny rule is
added only after each required flow passes independently.
\end{gate}

\begin{tabularx}{\linewidth}{@{}p{.22\linewidth}p{.31\linewidth}L@{}}
\toprule
\textbf{Attachment} & \textbf{What it buys} & \textbf{Cost or risk}\\
\midrule
Existing LAN & Fastest first hardware trial & Weakest destructive-install boundary; noisy policy\\
Dedicated interface & Existing host path stays untouched & Needs a second cable and switch port\\
VLAN trunk & Several isolated services over one cable & Bridge and recovery procedure are more complex\\
\bottomrule
\end{tabularx}

\proofpair{make homelab-workstation-verify
INSTANCE=<private-acceptance.json>}{Validates the recorded DHCP, AD DNS, boot
HTTP, time and cross-zone acceptance evidence without changing the network.}
\askbefore{Which physical port is the build port? Can an ordinary household
client see the installer? Can a build client reach a management interface? Is
there a tested local-console recovery path if the bridge is wrong?}

\section{Stage 5 --- Configure UniFi, one dependency at a time}

Create the small provisioning network at the expected device count, leaving
adjacent space available. Prefer simple human-facing gateway and service
addresses such as \texttt{.1}; the actual CIDR boundary remains technically
valid. Configure DHCP, AD DNS, time and internet update access before setting
network boot options. Add isolation rules last.

\begin{enumerate}
\item Create the VLAN and subnet. \textit{This gives build traffic a bounded
address space.} Verify one ordinary lease, gateway reachability and the exact
CIDR.
\item Set the network's DNS server to the AD DNS address. \textit{This enables
domain service discovery.} Verify the LDAP and Kerberos SRV records from a
client on that VLAN.
\item Permit time synchronization. \textit{Kerberos rejects sufficiently
skewed clocks.} Verify both client and DC report synchronized time.
\item Configure network boot for x86-64 UEFI, with the Controller address and
iPXE filename. \textit{This hands firmware to the versioned HTTP boot chain.}
Verify a test VM or laptop reaches the menu.
\item Permit only required update and service flows. \textit{This makes
installation useful without exposing management networks.} Test each allow
rule before the final deny.
\end{enumerate}

\failurebranch{If isolation breaks a client, inspect DHCP, DNS, time, then the
specific destination port. Do not disable the entire policy as a diagnostic
shortcut; enable or log one dependency at a time.}

\section{Stage 6 --- Mint one disposable pilot}

Connect the laptop through a known UEFI-PXE-capable USB-C Ethernet adapter.
Disable Secure Boot, select UEFI network boot, and choose the approved
workstation profile. The installer asks for the bespoke hostname, disk serial,
Windows/Arch allocation and default OS. It rejects a disk that cannot satisfy
both filesystem minimums.

\begin{center}
\begin{tikzpicture}[x=1cm,y=1cm,font=\scriptsize]
  \draw[wf pencil] (-5.7,-.55) rectangle (5.65,1.45);
  \foreach \x in {-5.05,-2.65,-.25,2.15,4.55}
    \draw[wf faint] (\x,-.55)--(\x,1.45);
  \node[align=center] at (-5.38,.45) {ESP\\1 GiB};
  \node[align=center] at (-3.85,.45) {Windows 11\\75\%\\min 160 GiB};
  \node[align=center] at (-1.45,.45) {recovery\\image-sized};
  \node[align=center] at (.95,.45) {Arch root\\25\%\\min 64 GiB};
  \node[align=center] at (3.80,.45) {free/derived\\alignment};
  \draw[wf pencil,<->] (-5.68,-.82)--(5.63,-.82);
  \node[below] at (0,-.83) {one GPT disk; exact sectors frozen before authorization};
\end{tikzpicture}
\wfcaption{The ratio is a policy default, never permission to violate either minimum.}
\end{center}

\begin{gate}{Destructive authorization}
The operator must compare the visible physical device with model, serial and
capacity, then type the complete disk serial. A profile selection, hostname or
\texttt{yes} is not authorization.
\end{gate}

\expectthis{Windows Setup completes first; Arch fills its declared partition;
both write independent UEFI entries; Windows becomes the default after a
five-second menu. The system does not join the domain until both native boots
pass.}
\proofpair{efibootmgr -v}{Windows and Arch entries both exist; the documented
default points to the five-second selector.}
\proofpair{lsblk -o NAME,SIZE,FSTYPE,PARTLABEL,UUID}{Observed partitions match
the frozen preflight plan within declared alignment tolerances.}

% Detailed dual-boot installation procedure.
% Included by main.tex after the pilot's disk plan has been frozen.

\section{Pilot installation --- measure, write, prove}

The profile is a proposed geometry, not permission to erase a disk. Print the
plan, then compare it with the laptop in front of you. Use sectors as the
authoritative unit; rounded GiB figures are only a human cross-check.

\begin{enumerate}
\item Photograph the firmware summary and the physical disk label if it is
accessible. Record firmware revision, disk model, serial and advertised
capacity.
\item Boot the approved PXE release in UEFI mode. Record the release version
and manifest digest shown by the installer.
\item Run the read-only preflight. Save \texttt{lsblk -b -o
NAME,SIZE,MODEL,SERIAL,TYPE,TRAN} and \texttt{wipefs -n} output.
\item Compare the discovered logical-sector size, last usable sector and disk
serial with the frozen plan. Differences are a stop condition.
\end{enumerate}

\askbefore{Is this the one authorized internal disk? Is the PXE or USB medium
absent from the target list? Does the disk have at least the Windows and Arch
minimums after the ESP, Microsoft reserved partition and recovery allowance?
Are there files on the existing disk that have not been independently
recovered and opened?}

\begin{tabularx}{\linewidth}{@{}p{.24\linewidth}p{.22\linewidth}p{.24\linewidth}L@{}}
\toprule
\textbf{Measurement} & \textbf{Frozen value} & \textbf{Observed value} & \textbf{Pass rule}\\
\midrule
Disk model and serial & \placeholder{plan} & \placeholder{preflight} & Exact serial; model consistent\\
Disk bytes & \placeholder{plan} & \placeholder{lsblk} & Exact byte count\\
Logical sector bytes & \placeholder{plan} & \placeholder{blockdev} & Exact match\\
First/last usable LBA & \placeholder{plan} & \placeholder{sgdisk} & Exact match\\
PXE release digest & \placeholder{release} & \placeholder{screen} & Exact SHA-256\\
\bottomrule
\end{tabularx}

\begin{gate}{Last reversible point}
Do not type the disk serial until the table passes. The next operation destroys
the selected partition table. No hostname, device name such as
\texttt{/dev/nvme0n1}, or general confirmation substitutes for the complete
serial.
\end{gate}

\subsection{Create and verify the GPT}

The automation creates, in order, one shared EFI System Partition (ESP), one
Microsoft Reserved partition (MSR), the Windows partition, Windows recovery,
and the Arch root partition. Start each resizable partition on the alignment
boundary emitted by the profile. Do not calculate boundaries by eye in a GUI.

\begin{center}
\begin{tikzpicture}[x=1cm,y=1cm,font=\scriptsize]
  \draw[wf pencil] (-5.75,-.20) rectangle (5.70,1.10);
  \draw[wf faint] (-5.22,-.20)--(-5.22,1.10);
  \draw[wf faint] (-4.93,-.20)--(-4.93,1.10);
  \draw[wf faint] (1.25,-.20)--(1.25,1.10);
  \draw[wf faint] (2.05,-.20)--(2.05,1.10);
  \node[rotate=90] at (-5.49,.45) {ESP};
  \node[rotate=90] at (-5.08,.45) {MSR};
  \node at (-1.84,.45) {Windows};
  \node[rotate=90] at (1.65,.45) {WinRE};
  \node at (3.87,.45) {Arch root};
  \draw[wf pencil,<->] (-5.75,-.47)--(-5.22,-.47);
  \node[below] at (-5.49,-.48) {1 GiB};
  \draw[wf pencil,<->] (-4.93,-.47)--(1.25,-.47);
  \node[below] at (-1.84,-.48) {profile-derived};
  \draw[wf pencil,<->] (2.05,-.47)--(5.70,-.47);
  \node[below] at (3.87,-.48) {profile-derived};
\end{tikzpicture}
\wfcaption{One GPT and one shared ESP; exact LBA boundaries come from the signed plan.}
\end{center}

\expectthis{A second read of the GPT reports the planned partition count,
type GUIDs, start sectors and end sectors. No filesystem signature exists
outside the partitions declared by the profile.}
\failurebranch{If any start sector, type GUID, serial or disk byte count differs,
stop before formatting. Save both observations, remove installer power, and
regenerate the plan for the hardware actually present.}

\subsection{Install Windows first}

Boot Windows Setup from the versioned PXE entry. Load only the driver bundle
whose hardware identifiers and digest were approved with the release. At the
disk chooser, select the pre-created partition labelled for Windows; do not
delete the table and do not allow Setup to consume the Arch partition.

\begin{enumerate}
\item Confirm Setup reports the planned Windows partition size within its
display-rounding tolerance. Ask: ``Which partition will Setup format, and what
two observed facts identify it?''
\item Install Windows 11 Pro. Firmware activation may occur later; activation
status is not an installation gate.
\item At the first desktop, install firmware and drivers from the laptop
vendor or Windows Update. Record each device that retains a warning icon.
\item Apply all update rounds, restarting and checking again until two
successive scans report no required update. Optional preview drivers remain
excluded unless a recorded hardware fault requires one.
\end{enumerate}

\proofpair{Get-Disk | Format-Table Number,FriendlyName,SerialNumber,Size}{The
system disk serial and byte count identify the authorized disk.}
\proofpair{Get-Partition | Sort-Object Offset}{Partition offsets and sizes
match the frozen plan; no unexpected data partition exists.}
\proofpair{Get-PnpDevice | Where-Object Status -ne 'OK'}{Returns no unexplained
device; record any intentional disabled device.}
\proofpair{Confirm-SecureBootUEFI}{Reports \texttt{False}, consistent with the
phase-1 boundary; an unsupported result is investigated, not assumed.}
\proofpair{manage-bde -status}{All workstation volumes report protection off;
device encryption must not have silently enabled itself.}

\begin{gate}{Windows recovery proof}
Open the Windows Recovery Environment once before installing Arch. Confirm the
keyboard works and that Startup Repair and Command Prompt are visible. Exit
without changing the disk. This proves the recovery partition is bootable
before another operating system changes UEFI entries.
\end{gate}

\subsection{Install Arch into its declared partition}

Return to the same PXE release and select Arch. Re-run disk identity preflight;
a successful Windows boot does not waive serial authorization. Format only the
Arch root partition. Mount the existing ESP at \texttt{/efi}; never format it.

\begin{enumerate}
\item Check time synchronization before package installation. Record clock
offset and time source; Kerberos will depend on this later.
\item Mount Arch root, then mount the existing FAT32 ESP. List
\texttt{/efi/EFI/Microsoft/Boot/bootmgfw.efi}; its presence is a stop-check
before installing a boot manager.
\item Install the base system, Intel microcode, NetworkManager, an LTS-capable
kernel path, and the approved graphics, Wi-Fi and firmware packages.
\item Generate \texttt{/etc/fstab} by UUID. Inspect every line: root must use
the Arch filesystem UUID and the ESP must use the shared FAT UUID.
\item Install \texttt{systemd-boot} without removing the Microsoft directory.
Create the Arch entry, set a five-second timeout, and set Windows
\texttt{auto-windows} as the default.
\end{enumerate}

\proofpair{findmnt --verify --verbose}{Reports no fstab parse or unreachable
source error.}
\proofpair{bootctl status}{Shows the shared ESP, an Arch entry and the automatic
Windows entry; default is Windows and timeout is five seconds.}
\proofpair{Verify the Windows EFI loader exists}{Confirm
\texttt{/efi/EFI/Microsoft/Boot/bootmgfw.efi} remains present after
the Arch boot manager is installed.}
\proofpair{lspci -k}{Each storage, graphics, Ethernet and wireless controller
has the expected kernel driver bound.}
\proofpair{systemctl --failed}{Returns no unexplained failed unit after a native
Arch boot.}

\failurebranch{If the automatic Windows entry is absent, do not hand-edit or
copy Microsoft binaries. Verify the ESP mount and the original Microsoft path.
If Windows files are missing, use Windows recovery media to repair them, then
rerun \texttt{bootctl status}.}

\subsection{Prove both native boot paths}

Power the laptop fully off between tests. A reboot alone may preserve firmware
or device state and conceal a cold-start defect.

\begin{enumerate}
\item Power on without touching a key. Measure from menu appearance until
Windows begins loading. Pass only if Windows is selected after five seconds.
\item Restart, select Arch, and measure time to a local login prompt. Record
network association time separately from operating-system boot time.
\item Use the firmware's one-time boot menu to start Windows Boot Manager
directly, then to start Linux Boot Manager directly.
\item Repeat once with Ethernet disconnected. Both local operating systems must
reach login independently of PXE, AD, DNS and network storage.
\end{enumerate}

\begin{tabularx}{\linewidth}{@{}p{.27\linewidth}p{.20\linewidth}p{.24\linewidth}L@{}}
\toprule
\textbf{Boot path} & \textbf{Expected} & \textbf{Measured} & \textbf{Evidence}\\
\midrule
Default menu & Windows, 5 s & & video / stopwatch\\
Selected Arch & local login & & journal boot ID\\
Firmware Windows entry & Windows login & & firmware photo\\
Firmware Linux entry & Arch login & & firmware photo\\
Both, cable absent & local login & & two boot IDs\\
\bottomrule
\end{tabularx}

\subsection{Recovery order}

Recovery preserves the other operating system before it improves convenience.
Do not repartition to repair a boot menu.

\begin{enumerate}
\item \textbf{Firmware first:} open the one-time menu. If either native entry
works, boot it and record \texttt{efibootmgr -v} or
\texttt{bcdedit /enum firmware}.
\item \textbf{Selector second:} from Arch media, mount the existing root and
ESP, confirm their UUIDs, chroot, then reinstall systemd-boot and regenerate
the Arch entry. Do not format either filesystem.
\item \textbf{Windows files third:} if the Microsoft EFI binary is absent or
Windows reports boot damage, use the matching Windows 11 recovery image and
\texttt{bcdboot} against the existing Windows and ESP volumes.
\item \textbf{Rollback last:} select the previous verified PXE release if a
current installer or driver bundle is suspect. A PXE rollback does not modify
an already installed workstation.
\end{enumerate}

\askbefore{Which layer failed: firmware entry, selector file, OS loader,
filesystem, or operating system? Which read-only observation proves that
classification? Can the repair be performed without formatting or moving a
partition? Where is the recovery evidence saved?}

\begin{gate}{Escalation boundary}
If neither filesystem is readable, the disk reports new media errors, or the
observed partition boundaries differ from the signed plan, stop. Image the disk
or replace the disposable pilot; do not use filesystem repair, partition
reconstruction, or installation retry as an exploratory diagnostic.
\end{gate}


\section{Stage 7 --- Join and test Windows}

Join Windows 11 Pro to \placeholder{identity-domain}, restart, and log in online once
with an ordinary domain test account. Test the daily administrator separately;
do not use the domain-wide administrator for ordinary desktop work.

\proofpair{whoami}{Returns the short domain and the intended user.}
\proofpair{whoami /groups}{The ordinary test user lacks local Administrators;
the approved workstation administrator has it.}
\proofpair{klist}{Kerberos tickets name the permanent realm.}
\proofpair{Resolve-DnsName -Type SRV \_ldap.\_tcp.\placeholder{identity-domain}}{The client
discovers the DC through AD DNS rather than a hosts-file entry.}

Disconnect the network and restart. The previously tested user must be able to
open the cached local profile. Reconnect, disable the account in AD, and prove
connected login is denied. Record that disconnected cached login remains a
known phase-1 limitation.

\section{Stage 8 --- Join and test Arch}

Configure SSSD against the same domain. Create local homes on first login;
never make the base home path depend on network storage. Grant sudo only
through an explicit AD workstation-administrator group.

\proofpair{id <test-user>}{Returns a stable domain UID, primary group and
expected memberships.}
\proofpair{kinit <test-user> \&\& klist}{The realm ticket succeeds while
online; destroy the ticket after recording the result.}
\proofpair{getent passwd <test-user>}{The NSS record and home path are present
without editing local \texttt{/etc/passwd}.}
\proofpair{sudo -l -U <test-user>}{The ordinary test user has no sudo rule.}

Repeat online login, offline cached login, connected account disable, DNS
failure and time-skew tests. A friendly desktop symptom is not evidence:
record command output and the exact expected boundary.

\% Integration point: place after Stage 8's local-rescue prerequisite and
\% before optional storage. This file deliberately contains no document wrapper.
\subsection{Join and prove Windows 11}

Start on a wired provisioning connection. Log in as the tested local rescue
administrator, set the permanent workstation name, and point the interface at
AD DNS. Do not join while the client is using public, gateway, or fallback
DNS.

\begin{shellblock}
Rename-Computer -NewName "<workstation-name>" -Restart
Get-DnsClientServerAddress -AddressFamily IPv4
Resolve-DnsName -Type SRV "\_ldap.\_tcp.dc.\_msdcs.<identity-domain>"
w32tm /query /source
Test-NetConnection <dc-fqdn> -Port 445
\end{shellblock}

\expectthis{The computer name is exact; every active domain-member interface
lists only approved AD DNS servers; the SRV answer names a DC; time has a
source; TCP 445 succeeds.}
\failurebranch{If ordinary names resolve but the SRV query fails, correct DHCP
or interface DNS. If the clock differs from the DC by five minutes or more,
correct time before touching the domain. Do not work around either fault with
a hosts-file entry.}

Join from \textit{Settings > System > About > Domain or workgroup > Change},
select \textit{Domain}, enter \placeholder{identity-domain}, and supply the
temporary delegated join credential. A domain administrator is not required
for routine joins. Restart when Windows asks.

\begin{shellblock}
Get-CimInstance Win32\_ComputerSystem |
  Select-Object Name,PartOfDomain,Domain
nltest /dsgetdc:<identity-domain>
Test-ComputerSecureChannel -Verbose
\end{shellblock}

\expectthis{\texttt{PartOfDomain} is true, \texttt{Domain} is exact,
\texttt{nltest} discovers a DC through DNS, and the secure channel reports
\texttt{True}. Record the DC name and test time.}

At the sign-in screen choose \textit{Other user}. Sign in once as the standard
test identity using \texttt{DOMAIN\textbackslash user} or the full user
principal name. Open a non-elevated PowerShell session:

\begin{shellblock}
whoami
whoami /groups
echo \$env:USERPROFILE
net localgroup Administrators
\end{shellblock}

\proofpair{The reported identity is the standard domain test user}{Its profile
is local, its SID is present, and it is absent from the local Administrators
group.}
\proofpair{Sign in as the approved daily administrator}{A deliberate elevation
prompt succeeds, but the reserved domain-administrator identity has not been
used for daily work.}
\proofpair{Disconnect every network interface and restart}{The previously used
domain identity signs in from its cached credential; network resources fail
without delaying or preventing the local desktop.}
\proofpair{Reconnect, disable the standard test identity, then attempt a fresh
online sign-in}{Connected authentication is denied. Record separately that an
already cached offline credential cannot be revoked in phase 1.}

\askbefore{Did the join use a delegated credential? Can local rescue still sign
in with the DC powered off? Is the standard identity non-administrative? Did
the test measure DNS answer, clock source, secure-channel state, and profile
path rather than infer success from a desktop appearing?}

\subsection{Join and prove Arch Linux}

Finish a full system update before the join; Arch does not support partial
upgrades. Install the declared identity client as one transaction:

\begin{shellblock}
sudo pacman -Syu --needed sssd samba krb5 bind
resolvectl status
timedatectl show -p NTPSynchronized --value
host -t SRV \_ldap.\_tcp.dc.\_msdcs.<identity-domain>
\end{shellblock}

\expectthis{AD DNS is active, time reports \texttt{yes}, the LDAP SRV answer
names a DC, and the packages came from enabled official Arch repositories.}
\failurebranch{Stop if discovery fails, time is not synchronized, or Network
Manager can replace AD DNS with a public resolver. A successful ping is not an
identity-plane test.}

Create a one-use delegated join identity or obtain a time-limited credential
from the encrypted secret channel. Generate the Samba configuration from
private values with \texttt{security = ads}, the exact
\texttt{workgroup = <netbios-name>} and
\texttt{realm = <kerberos-realm>}. Set the Kerberos method to
\texttt{secrets and keytab}. Generate the root-owned, mode-0600 SSSD
configuration with:

\begin{shellblock}
[sssd]
services = nss, pam
config\_file\_version = 2
domains = <identity-domain>

[domain/<identity-domain>]
id\_provider = ad
auth\_provider = ad
access\_provider = ad
ad\_domain = <identity-domain>
krb5\_realm = <kerberos-realm>
cache\_credentials = true
use\_fully\_qualified\_names = true
fallback\_homedir = /home/\%u
default\_shell = /bin/bash
ldap\_id\_mapping = true
\end{shellblock}

Review both rendered files before joining. Join interactively; never place the
password in a command, shell history, answer file, or repository.

\begin{shellblock}
sudo net ads join -U '<delegated-join-user>'
sudo net ads testjoin
sudo klist -k /etc/krb5.keytab
sudo test "\$(stat -c \%a /etc/sssd/sssd.conf)" = 600
sudo systemctl enable --now sssd
\end{shellblock}

\expectthis{\texttt{net ads testjoin} reports success, the machine keytab
contains the expected host principals, \texttt{/etc/sssd/sssd.conf} is
root-owned mode 0600, and \texttt{sssd.service} is active. Revoke or expire
the join credential now.}

Keep the naming rule explicit. Either retain fully qualified login names
(\texttt{user@domain}) or set a reviewed, site-wide short-name rule; never
allow different domains or local accounts to collapse onto the same visible
name. Confirm that \texttt{sss} appears on the \texttt{passwd}, \texttt{group},
and \texttt{shadow} lines in \texttt{/etc/nsswitch.conf}. Then verify lookups:

\begin{shellblock}
getent passwd '<standard-user>@<identity-domain>'
id '<standard-user>@<identity-domain>'
getent group '<approved-admin-group>@<identity-domain>'
sssctl user-checks '<standard-user>@<identity-domain>' -s login
\end{shellblock}

\expectthis{The user resolves once, has a stable numeric UID and GID, is absent
from the approved administrator group, and SSSD authorizes the login service.
Record the numeric IDs; optional SMB files must present compatible ownership.}

Local home creation is a PAM action, not a network-home mount. Before editing,
save the package-owned file's digest and a private evidence copy. Add this line
to the session stack used by the workstation login path, normally
\texttt{/etc/pam.d/system-login}, after the other required session setup:

\begin{shellblock}
session optional pam\_mkhomedir.so skel=/etc/skel umask=0077
\end{shellblock}

Review the complete PAM stack before rebooting. Do not add
\texttt{pam\_mkhomedir} to authentication or account sections, and do not
replace package-managed includes. Test first from a second TTY while the local
rescue session remains open.

\begin{shellblock}
getent passwd '<standard-user>@<identity-domain>'
findmnt --target '/home/<resolved-home>'
stat -c '\%U \%G \%a \%n' '/home/<resolved-home>'
sudo -l -U '<standard-user>@<identity-domain>'
\end{shellblock}

\proofpair{First online login}{Creates exactly the home path returned by
\texttt{getent}; it is a local filesystem path, owned by the resolved UID/GID,
and begins with restrictive permissions.}
\proofpair{Standard-user authorization}{The user receives no workstation
administration through \texttt{sudo}.}
\proofpair{Approved administrator authorization}{Only the reviewed domain
group receives the intended narrow \texttt{sudo} policy; a real elevation is
tested and logged.}
\proofpair{Offline restart after one successful online login}{The cached
identity signs in, the same numeric UID and GID are returned, and the existing
local home opens even when DNS, LDAP, Kerberos, and SMB are unavailable.}
\proofpair{Local rescue with the DC unavailable}{Login and \texttt{sudo}
succeed without SSSD. If they do not, restore the PAM/NSS configuration from
the still-open rescue session.}

\begin{gate}{UID, time, and storage proof}
Record UID, primary GID, supplementary groups, client time, DC time, and the
SMB file owner observed from both operating systems. A name that merely looks
correct is not proof. NFS remains disabled until numeric-ID mapping is designed
and tested; phase 1 uses SMB and a local home.
\end{gate}

\askbefore{What exact login spelling is canonical? Which AD group alone grants
local administration? Does the home exist and open with storage powered off?
Do UID and GID remain unchanged across two cold boots? Is clock offset below
the Kerberos limit? Does disabling an account deny a new connected login while
the documented offline-cache limitation remains understood?}


\section{Stage 9 --- Add optional SMB storage}

Integrate the NAS with AD before touching a workstation. Give each person a
private share, or a folder whose ACL grants that AD person and the storage
administrators access. Do not store a second copy of the person's password.
The workstation presents the current Kerberos identity to SMB.

\begin{gate}{The local profile is authoritative}
Windows must reach its desktop and Arch must reach its local home without the
NAS. Network storage is a convenience mounted after login. Never redirect
\texttt{Desktop}, \texttt{Documents}, \texttt{\%USERPROFILE\%}, or
\texttt{/home} during phase 1.
\end{gate}

\subsection*{Prove the server before configuring a client}

\begin{enumerate}
\item From AD DNS, resolve \placeholder{storage-host}. Record the returned
address and resolver. \textit{This catches split-DNS and stale-record errors.}
\item With an ordinary test identity, list the intended share. Confirm that a
different ordinary identity is denied. \textit{This proves the server ACL
instead of relying on a drive letter or desktop icon.}
\item Create a uniquely named test file, read it back, remove it, and record
the server-side owner. \textit{This detects guest fallback and identity-mapping
errors.}
\item Compare client, domain-controller, and NAS time. Keep observed skew below
the site's Kerberos limit. \textit{SMB may look like an ACL failure when the
ticket is actually too old or the clocks disagree.}
\end{enumerate}

\failurebranch{If the owner is guest, an unexpected numeric identity, or
another household user, stop. Correct the NAS domain join and ACL model before
automating any workstation mount. Do not compensate with a shared credential
or a permissive ACL.}

\subsection*{Windows: map after desktop login}

\begin{enumerate}
\item Log in online once as the ordinary domain test user. Run
\texttt{klist}; obtain a ticket before continuing.
\item In File Explorer, map the approved UNC path to the documented drive
letter. Select reconnect, but do not select alternate credentials and do not
save a password in Credential Manager.
\item Put the mapping action in a per-user logon task with an eight-second
timeout and no ``wait for network'' prerequisite. The task may report failure;
it may not hold the desktop open. Record the task name and exported,
secret-free definition.
\item Sign out and back in with the NAS reachable. Create a file, disconnect
the mapping, reconnect it, and read the same bytes. Record the path, SHA-256
digest, elapsed login time, and effective AD owner.
\end{enumerate}

\proofpair{Get-SmbMapping}{The approved drive points to the approved UNC path,
uses the signed-in domain identity, and reports no guest mapping.}
\proofpair{cmd /c ``net use''}{No obsolete drive letter, alternate NAS name,
or unexpected cached connection is present.}

\subsection*{Arch: mount in the user session}

Install the repository-declared SMB and desktop virtual-filesystem packages.
Use the desktop's user-session SMB mount, such as GVfs, so no network
filesystem participates in boot or local-home creation.

\begin{enumerate}
\item Log in online as the ordinary domain test user. Run \texttt{klist}; fix
SSSD or Kerberos first if no current ticket exists.
\item Open \texttt{smb://\placeholder{storage-host}/\placeholder{user-share}}
from the desktop file manager. Choose the current domain identity and decline
permanent password storage.
\item If automatic connection is wanted, create a user-session service that
runs the mount after the graphical session starts, has an eight-second
timeout, and is not required by the session target. Never add the SMB path to
\texttt{/etc/fstab}, the boot target, or the local-home unit.
\item Create, hash, unmount, remount, and read the same test file. Record the
Kerberos principal and effective server-side owner.
\end{enumerate}

\proofpair{gio mount -l}{The approved SMB location appears only when reachable;
the local home remains an ordinary local filesystem.}
\proofpair{systemctl --user --failed}{A failed optional mount may be visible,
but no login, home, or graphical-session unit depends on it.}

\subsection*{Run the three-state failure drill on both systems}

\begin{tabularx}{\linewidth}{@{}p{.17\linewidth}p{.31\linewidth}L@{}}
\toprule
\textbf{State} & \textbf{Action and measurement} & \textbf{Pass boundary}\\
\midrule
Available & Login, create and hash a file; unmount and remount & Same bytes,
same AD owner, no password prompt\\
Unreachable & Block or disconnect only the NAS; cold-login with a stopwatch &
Local desktop/home opens within the declared limit; bounded warning only\\
Denied & Keep the host reachable but remove the test user's share access &
Explicit access-denied result; local login works; no guest or alternate user\\
\bottomrule
\end{tabularx}

\askbefore{Did each test begin from a cold sign-in rather than an already
connected session? Was the NAS itself unreachable in the second test, but
reachable in the third? Did either client silently reuse an old connection or
credential?}

\failurebranch{If login waits on storage, remove the mapping from every boot,
profile, shell-startup and required-session dependency, then repeat all three
states. If access denied becomes access granted after a prompt, remove the
saved credential and inspect the server ACL.}

\section{Stage 10 --- Add restricted managed Wi-Fi}

Create a dedicated managed-workstation SSID and VLAN. Prefer one PPSK per
workstation so one lost laptop can be removed without changing every machine.
Retrieve the selected key from the private secret store only while minting the
machine. Phase-1 disks are unencrypted; assume physical access can recover any
stored wireless credential.

\begin{center}
\begin{tikzpicture}[x=1cm,y=1cm,font=\scriptsize]
  \draw[wf pencil,rounded corners=1pt] (-5.6,-.10) rectangle (-3.35,1.22);
  \draw[wf faint] (-5.42,.08) rectangle (-3.53,1.04);
  \node[align=center] at (-4.48,.56) {managed\\workstation};
  \draw[wf flow] (-3.35,.56)--(-1.63,.56) node[midway,above]{PPSK};
  \draw[wf pencil] (-1.63,.18)--(-1.15,.56)--(-1.63,.94);
  \draw[wf pencil] (-1.35,.05)--(-.71,.56)--(-1.35,1.07);
  \draw[wf pencil] (-1.03,-.09)--(-.20,.56)--(-1.03,1.21);
  \node[below] at (-.92,-.13) {restricted SSID};
  \draw[wf flow] (-.20,.56)--(1.10,.56);
  \draw[wf pencil] (1.10,-.10) rectangle (2.65,1.22);
  \node[align=center] at (1.88,.56) {AD/DNS\\updates\\approved SMB};
  \draw[wf deny] (3.12,-.10)--(3.12,1.22);
  \draw[wf pencil] (2.91,1.06)--(3.33,.06);
  \draw[wf pencil] (2.91,.06)--(3.33,1.06);
  \draw[wf faint] (3.72,-.10) rectangle (5.58,1.22);
  \node[align=center] at (4.65,.56) {management\\household peers};
\end{tikzpicture}
\wfcaption{A stored credential reaches only the services a managed workstation needs.}
\end{center}

\subsection*{Prove UniFi policy before storing a key}

Using a temporary test client, verify association, DHCP lease, gateway, AD
DNS, time, domain ports, approved SMB, and both operating systems' update
destinations. Then verify denials to gateway administration, Controller
administration, infrastructure management, household-client, peer-client, and
IoT destinations. Save a small source/destination/port/result table.

\failurebranch{If an allow fails, inspect association, DHCP, DNS, time, then
the single destination and port. If a deny fails, quarantine the SSID. Do not
store its credential on a workstation until both directions match policy.}

\subsection*{Windows: install one machine profile}

\begin{enumerate}
\item Create a WLAN profile for only the managed SSID and its assigned PPSK.
Import it for all users from the local console. Do not add household,
management, or provisioning SSIDs.
\item Delete the plaintext import file immediately after the profile appears.
Check Git status, the build log, answer files and shell history for the SSID
key; store only the key's secret-manager reference in evidence.
\item Disconnect Ethernet, restart, and log in with a previously cached
ordinary account. Confirm automatic association and the expected managed-VLAN
lease.
\item Forget or disable the profile, restart, and prove that Windows does not
fall back to another known private network. Restore only the approved profile.
\end{enumerate}

\proofpair{netsh wlan show profiles}{Exactly the approved preload set appears;
captured output contains no key material.}
\proofpair{Get-NetConnectionProfile}{The active wireless interface has the
expected network category and no second active route into a private zone.}

\subsection*{Arch: install one root-owned NetworkManager profile}

\begin{enumerate}
\item Create one system connection for the managed SSID with NetworkManager,
binding it to the workstation PPSK and enabling autoconnect. Avoid typing the
key in a command that will be retained in shell history.
\item Confirm its file under
\texttt{/etc/NetworkManager/system-connections/} is owned by root, mode 600,
and excluded from logs, artifacts, backups without secret encryption, and Git.
\item Disconnect Ethernet, restart NetworkManager or reboot, and log in with a
previously cached ordinary account. Confirm association, expected lease, AD
DNS, time, and the default route.
\item Set the connection down and prove that no household or management SSID
autoconnects. Bring back only the approved profile.
\end{enumerate}

\proofpair{nmcli -f NAME,TYPE,AUTOCONNECT connection show}{Only approved
wireless profiles autoconnect; do not include secret fields in evidence.}
\proofpair{stat -c '\%U \%G \%a \%n'
/etc/NetworkManager/system-connections/*}{Every stored system profile is
root-owned and mode 600.}

\subsection*{Wireless-only acceptance}

Repeat AD DNS SRV lookup, clock check, ordinary online login, administrator
authorization, OS update check and optional-SMB three-state drill with Ethernet
physically disconnected. Measure association-to-lease and login time. Scan
each forbidden zone from only its documented small CIDR; do not substitute a
wide private-range scan.

\askbefore{Does revoking this workstation's PPSK leave every other workstation
connected? Does the revoked machine fail to associate after its cached radio
state is cleared? Can it reach all required update endpoints but no management
page or peer workstation?}

\failurebranch{If wireless onboarding fails, continue over the deployment
cable. Never inject the household Wi-Fi key or broaden the VLAN to make the
test pass. If a required vendor update cannot be bounded yet, record that
exception and keep the workstation a pilot.}


\section{Stage 11 --- Publish and migrate}

Publishing a PXE release uploads to a new immutable version directory, verifies
the remote bytes with a checksum comparison, and atomically moves
\texttt{current}. Rollback selects a previous verified release; it never edits
that release.

\begin{shellblock}
make homelab-pxe-publish RELEASE=<local-release> \textbackslash
  DESTINATION=<host:/absolute/root>
make homelab-pxe-rollback TARGET=<target> VERSION=<previous-version> \textbackslash
  DESTINATION=<host:/absolute/root>
\end{shellblock}

The permanent physical Controller later joins as \texttt{permanent\_dc\_fqdn}.
Validate AD/DNS replication. Move the boot alias and UniFi boot settings, then
demote \texttt{bootstrap-dc}. Workstations are not rebuilt because their
permanent domain, realm and discovery mechanism do not change.

\begin{telosfacts}[0.27]
\telosfact{Future VM candidate}{Samba AD/DNS can move behind the same names after an additional DC replicates and roles transfer.}
\telosfact{Future VM candidate}{A deployment credential broker can issue short-lived joins without placing reusable credentials in artifacts.}
\telosfact{Future VM candidate}{Windows deployment helpers, certificate services, audit collection and management services can be isolated later.}
\telosfact{No rebuild trigger}{A service move is acceptable only when an existing Windows and Arch client passes discovery and login unchanged.}
\end{telosfacts}

\clearpage
\section{Final acceptance worksheet}

Complete this only after the procedure. Worksheets belong at the back so they
support the instruction rather than interrupt its flow.

\begin{tabularx}{\linewidth}{@{}p{.31\linewidth}p{.20\linewidth}L@{}}
\toprule
\textbf{Claim} & \textbf{Pass / fail} & \textbf{Evidence location or measured value}\\
\midrule
One DHCP authority observed & & \\
AD LDAP and Kerberos SRV records resolve & & \\
Controller and client clocks synchronized & & \\
PXE release version and manifest verified & & \\
Disk serial matched before first write & & \\
Windows/Arch allocation matches plan & & \\
Windows is default after five seconds & & \\
Ordinary user logs in on Windows & & \\
Ordinary user logs in on Arch & & \\
Ordinary user lacks admin and sudo & & \\
Workstation administrator has intended rights & & \\
Both systems permit tested offline login & & \\
Connected disabled account is denied & & \\
SMB unavailable does not block login & & \\
No plaintext secret appears in Git or logs & & \\
\bottomrule
\end{tabularx}

\medskip
\noindent Workstation hostname:\ \rule{.38\linewidth}{.35pt}\hfill
Disk serial:\ \rule{.30\linewidth}{.35pt}

\medskip
\noindent PXE release:\ \rule{.30\linewidth}{.35pt}\hfill
Operator/date:\ \rule{.35\linewidth}{.35pt}

\begin{gate}{Release decision}
Every row passes with evidence, or the workstation remains a disposable pilot.
Record exceptions explicitly; familiarity, a successful desktop login, and
``it worked once'' are not acceptance evidence.
\end{gate}

\subsection*{Exception and retest record}
\noindent Record each failed gate and its evidence; strike unused rows.

\renewcommand{\arraystretch}{1.65}
\begin{tabularx}{\linewidth}{@{}p{.08\linewidth}p{.34\linewidth}p{.35\linewidth}L@{}}
\toprule
\textbf{Gate} & \textbf{Failure and measurement} & \textbf{Correction and retest evidence} & \textbf{Result}\\
\midrule
 & & & \\
\cmidrule(lr){1-4}
 & & & \\
\bottomrule
\end{tabularx}
\renewcommand{\arraystretch}{1}

\smallskip
\noindent Evidence packet location/checksum:\ \rule{.60\linewidth}{.35pt}

\smallskip
\noindent Reviewer/date:\ \rule{.24\linewidth}{.35pt}\hfill
\textbf{Disposition:}\quad
\textbf{release}\quad/\quad\textbf{pilot only}\quad/\quad\textbf{rebuild}

\end{document}
